\chapter{Appendix}
\label{chap:appendix}

\section{SQL Code}
\noindent
All SQL queries used throughout Chapter~\ref{chap:dataset_understanding} are presented below. Each query is placed within a figure environment for clarity and reference.

\begin{figure}[H]
\centering
\begin{minipage}{\textwidth}
\begin{verbatim}
-- Figure A.1: Counting rows/patients in the 'patients' table
SELECT
    COUNT(*) AS total_rows_in_patients,
    COUNT(DISTINCT subject_id) AS num_unique_patients
FROM mimic.patients;
\end{verbatim}
\end{minipage}
\caption{SQL code to retrieve the total number of rows and distinct patients in the \texttt{patients} table.}
\label{fig:sql1_patient_rows}
\end{figure}

\begin{figure}[H]
\centering
\begin{minipage}{\textwidth}
\begin{verbatim}
-- Figure A.2: Counting rows/admissions in the 'admissions' table
SELECT
    COUNT(*) AS total_rows_in_admissions,
    COUNT(DISTINCT hadm_id) AS num_unique_admissions
FROM mimic.admissions;
\end{verbatim}
\end{minipage}
\caption{SQL code to retrieve the total number of rows and distinct admissions in the \texttt{admissions} table.}
\label{fig:sql2_admission_rows}
\end{figure}

\begin{figure}[H]
\centering
\begin{minipage}{\textwidth}
\begin{verbatim}
-- Figure A.3: Identifying multiple admissions per patient
SELECT subject_id, COUNT(*) AS num_admissions
FROM mimic.admissions
GROUP BY subject_id
HAVING COUNT(*) > 1
ORDER BY num_admissions DESC
LIMIT 5;
\end{verbatim}
\end{minipage}
\caption{SQL code to find patients with more than one admission, illustrating recurrent hospital stays.}
\label{fig:sql3_multiple_admissions}
\end{figure}

\begin{figure}[H]
\centering
\begin{minipage}{\textwidth}
\begin{verbatim}
-- Figure A.4: Counting total discharge notes
SELECT COUNT(*) AS total_discharge_notes
FROM mimic.discharge;
\end{verbatim}
\end{minipage}
\caption{SQL code to count the total number of discharge summaries in the \texttt{discharge} table.}
\label{fig:sql4_discharge_notes}
\end{figure}

\begin{figure}[H]
\centering
\begin{minipage}{\textwidth}
\begin{verbatim}
-- Figure A.5: Admissions with and without discharge notes
SELECT
    (SELECT COUNT(DISTINCT hadm_id) FROM mimic.admissions) AS total_admissions,
    (SELECT COUNT(DISTINCT hadm_id) FROM mimic.discharge)  AS hadm_with_discharge,
    (SELECT COUNT(DISTINCT hadm_id) FROM mimic.admissions
     WHERE hadm_id NOT IN (
         SELECT DISTINCT hadm_id FROM mimic.discharge
     )
    ) AS hadm_without_discharge;
\end{verbatim}
\end{minipage}
\caption{SQL code comparing the number of admissions with discharge notes to those without.}
\label{fig:sql5_hadm_no_discharge}
\end{figure}

\begin{figure}[H]
\centering
\begin{minipage}{\textwidth}
\begin{verbatim}
-- Figure A.6: Word count distribution of discharge summaries
SELECT
    PERCENTILE_CONT(0.5)
        WITHIN GROUP (ORDER BY (LENGTH(text) - LENGTH(REPLACE(text, ' ', '')) + 1))
      AS median_word_count,
    PERCENTILE_CONT(0.9)
      WITHIN GROUP (ORDER BY (LENGTH(text) - LENGTH(REPLACE(text, ' ', '')) + 1))
    AS p90_word_count,
    PERCENTILE_CONT(0.99)
    WITHIN GROUP (ORDER BY (LENGTH(text) - LENGTH(REPLACE(text, ' ', '')) + 1))
    AS p99_word_count
FROM mimic.discharge;
\end{verbatim}
\end{minipage}
\caption{SQL code for computing word-count percentiles of the \texttt{discharge} table.}
\label{fig:sql6_note_wordcount}
\end{figure}

\begin{figure}[H]
\centering
\begin{minipage}{\textwidth}
\begin{verbatim}
-- Figure A.7: Comparing total code assignments in ICD-9 vs. ICD-10
SELECT
    icd_version,
    COUNT(*) AS code_assignments
FROM mimic.diagnoses_icd
GROUP BY icd_version
ORDER BY code_assignments DESC;
\end{verbatim}
\end{minipage}
\caption{SQL code to assess the distribution of ICD-9 and ICD-10 usage in the dataset.}
\label{fig:sql7_icd_version_counts}
\end{figure}

\begin{figure}[H]
\centering
\begin{minipage}{\textwidth}
\begin{verbatim}
-- Figure A.8: Retrieving top ICD-10 codes by frequency
SELECT icd_code, COUNT(*) AS freq
FROM mimic.diagnoses_icd
WHERE icd_version = 10
GROUP BY icd_code
ORDER BY freq DESC
LIMIT 10;
\end{verbatim}
\end{minipage}
\caption{SQL code used to find the 10 most common ICD-10 codes in MIMIC-IV.}
\label{fig:sql8_top_icd10_codes}
\end{figure}

\begin{figure}[H]
\centering
\begin{minipage}{\textwidth}
\begin{verbatim}
-- Figure A.9: Counting unique ICD-10 codes
SELECT icd_version,
       COUNT(DISTINCT icd_code) AS num_unique_codes
FROM mimic.diagnoses_icd
GROUP BY icd_version;
\end{verbatim}
\end{minipage}
\caption{SQL code showing the total number of unique ICD codes by ICD version.}
\label{fig:sql9_unique_codes}
\end{figure}

\begin{figure}[H]
\centering
\begin{minipage}{\textwidth}
\begin{verbatim}
-- Figure A.10: Identifying ICD-10 codes with fewer than five occurrences
WITH code_counts AS (
    SELECT icd_code, COUNT(*) AS freq
    FROM mimic.diagnoses_icd
    WHERE icd_version = 10
    GROUP BY icd_code
)
SELECT COUNT(*) AS codes_under_5_occurrences
FROM code_counts
WHERE freq < 5;
\end{verbatim}
\end{minipage}
\caption{SQL code to quantify rare ICD-10 codes that appear fewer than five times.}
\label{fig:sql10_rare_codes}
\end{figure}

\begin{figure}[H]
\centering
\begin{minipage}{\textwidth}
\begin{verbatim}
-- Figure A.11: Calculating average ICD-10 codes per admission
WITH code_counts AS (
    SELECT hadm_id,
           COUNT(DISTINCT icd_code) AS code_count
    FROM mimic.diagnoses_icd
    WHERE icd_version = 10
    GROUP BY hadm_id
)
SELECT AVG(code_count) AS avg_icd10_codes_per_adm
FROM code_counts;
\end{verbatim}
\end{minipage}
\caption{SQL code to compute the average number of ICD-10 codes per admission.}
\label{fig:sql11_avg_codes_per_admission}
\end{figure}

% \begin{figure}[H]
% \centering
% \begin{minipage}{\textwidth}
% \begin{verbatim}
% -- Figure A.12: (Optional) Example of note length vs. code count
% -- This type of query can gather correlation data,
% -- though not shown in detail here.
% -- SELECT ...
% \end{verbatim}
% \end{minipage}
% \caption{An example placeholder for more specialized correlation queries (omitted details).}
% \label{fig:sql12_note_length_vs_codes}
% \end{figure}

\begin{figure}[H]
\centering
\begin{minipage}{\textwidth}
\begin{verbatim}
-- Figure A.12: Merging discharge summaries with associated ICD-10 codes
DROP TABLE IF EXISTS mimic.full_notes_with_codes;

CREATE TABLE mimic.full_notes_with_codes AS
SELECT
    d.subject_id,
    d.hadm_id,
    d.text AS notes,
    STRING_AGG(di.icd_code, ',') AS icd_codes
FROM mimic.discharge AS d
JOIN mimic.admissions AS a
  ON d.subject_id = a.subject_id
    AND d.hadm_id = a.hadm_id
JOIN mimic.patients AS p
  ON p.subject_id = d.subject_id
JOIN mimic.diagnoses_icd AS di
  ON di.subject_id = d.subject_id
    AND di.hadm_id = d.hadm_id
WHERE di.icd_version = 10  -- focusing on ICD-10-CM only
GROUP BY
    d.subject_id,
    d.hadm_id,
    d.text;
\end{verbatim}
\end{minipage}
\caption{SQL code to create a consolidated table that pairs each discharge summary with its ICD-10 codes.}
\label{fig:sql12_merge_notes_codes}
\end{figure}

\chapter{MinHash Pseudo-Code}
\label{app:minhash}

\begin{figure}[H]
\centering
\begin{verbatim}
Algorithm 1: MinHash for Similar Document Detection

Input:
    - A collection of textual discharge summaries (D)
    - Number of permutations (num_perm)
    - Similarity threshold (t)

Output:
    - Set of document pairs whose Jaccard similarity exceeds t

Procedure:
1. Preprocess each document by:
   a) Converting text to lowercase
   b) Removing punctuation and non-alphanumeric characters
   c) Eliminating stopwords
2. Initialize an LSH index with threshold t and num_perm permutations.
3. For each document d in D:
   a) Create a MinHash object m with num_perm permutations
   b) For each token in d, update m using m.update(token.encode('utf8'))
   c) Insert the MinHash signature m into the LSH index
4. For each document d in D:
   a) Query the LSH index with d’s signature to obtain similar documents
   b) Record document pairs that exceed the threshold
5. Return all identified pairs of documents with similarity above t
\end{verbatim}
\caption{Pseudocode for the MinHash Similarity Detection Process}
\end{figure}


% %%%%%%%%%%%%%%%%%%%%%%%%%%%%%%%%%%%%%%%%%

\chapter{Appendix: Pseudocode}

\section{Ontology-Guided Synthetic Note Generation}
\begin{figure}[H]
\centering
\caption{Algorithm 1: Ontology-Guided Synthetic Note Generation}
\label{fig:algo1}
\begin{verbatim}
Input:
  r: Rare ICD code
  D: Disease cluster from ontology
  L: Language model with retrieval augmentation

Output:
  S: Synthetic discharge summary

Step 1:
  context_docs = RetrieveRelevantText(D, knowledge_source)
Step 2:
  structured_prompt = GenerateStructuredPrompt(D, context_docs)
  S_draft = L.generate(structured_prompt)
Step 3:
  for iteration in 1..MaxIter:
    feedback = CriticEvaluate(S_draft)
    if feedback indicates no major issues:
      break
    S_draft = L.refine(S_draft, feedback)
Step 4:
  if Validate(S_draft):
    S = S_draft
  else:
    S = None
\end{verbatim}
\end{figure}

\section{Hybrid Training with Rare-Aware Strategies}
\begin{figure}[H]
\centering
\caption{Algorithm 2: Hybrid Training with Rare-Aware Strategies}
\label{fig:algo2}
\begin{verbatim}
Input:
  RealData = {(x_i, Y_i)}
  SyntheticData = {(x_j*, Y_j*)}
  Model parameters θ

Output:
  Trained model parameters θ

Step 1:
  TrainSet = RealData ∪ SyntheticData
  Initialize θ

Step 2:
  for epoch in 1..max_epochs:
    for batch B in WeightedSampleBatches(TrainSet):
      (X_B, Y_B) = PrepareBatch(B)
      if epoch >= epoch_focal_start:
        loss = FocalLoss(θ, X_B, Y_B) + BCE(θ, X_B, Y_B)
      else:
        loss = WeightedBCE(θ, X_B, Y_B)
      θ = θ - learning_rate * ∇_θ(loss)

Step 3:
  Evaluate on validation set and save best θ based on macro-F1.
\end{verbatim}
\end{figure}